\section{Fixed bounded formulas and Boolean reflection}
\label{sec:reflection}

The arithmetic compiler needs a uniform way to turn a finite verification
condition into quantifier-free equations.  The phrase ``bounded formula'' is
potentially misleading: a bounded quantifier is still a quantifier in the
syntax, whereas the induction scheme of \(\PRAAx\) accepts only formulas
recognized by \(\Fm.\qf\).  SCR therefore compiles bounded logic to
primitive-recursive data and proves reflection separately.

\subsection{Reflected atoms}

A reflected atom over \(n\) variables consists of a formula \(A\), a bit term
\(b_A\), a proof that \(b_A\) is Boolean, and derivations connecting its two
values with \(A\):
\[
\begin{aligned}
   &\turnstile b_A=0\;\vee\;b_A=1,\\
   &b_A=1\turnstile A,\\
   &A\turnstile b_A=1.
\end{aligned}
\]
Equality, order, divisibility with a supplied quotient, and all polynomial
identities used by the packet layer have such characteristic terms in the
checked arithmetic library or can be constructed from them.  Negation is
represented by complement, conjunction by multiplication of bits, and
disjunction by the sign of their sum.  Each combinator carries object proofs.

This is intentionally a \emph{specialization-time} interface.  A Lean
structure may hold \((A,b_A,\pi_A)\), and a recursive Lean function may combine
these structures.  It emits a fixed term and a fixed derivation.  No natural
number inside iPRA is asked to interpret arbitrary syntax.

\subsection{Bounded existential search}

Let \(b(i,q)\) be a Boolean term and put
\[
  z(i,q)=1-b(i,q),
\]
using the repository's truncated subtraction convention.  A witness with
\(b=1\) is a zero of \(z\).  Define
\[
   \mu(N,q)=\mathsf{Search.mu}(z,N,q).
\]
The checked theorem \code{search\_spec} gives
\[
\begin{aligned}
  &\mu(N,q)\le N,\\
  &\mu(N,q)<N\to z(\mu(N,q),q)=0,\\
  &i<\mu(N,q)\to z(i,q)\ne 0.
\end{aligned}
\]
Thus the quantifier-free formula
\[
   E_b(N,q)\;:\!\Longleftrightarrow\;
   \mu(N,q)<N\;\wedge\;b(\mu(N,q),q)=1
\]
serves as the reflected bounded existential.  From \(E_b\) one introduces the
ordinary existential with witness \(\mu\); from any ordinary witness
\(i<N\) with \(b(i,q)=1\), \code{search\_found} proves \(E_b\).  The proof of
existential reflection therefore uses the existing qf induction hidden inside
\code{Search}, not induction on a quantified formula.

\begin{theorem}[Bounded-existential reflection]
For every fixed reflected atom family \(A(i,q)\), there is a fixed
quantifier-free formula \(E_A(N,q)\) and iPRA derivations
\[
  E_A(N,q)\to\exists i\,(i<N\wedge A(i,q))
  \qquad	ext{and}\qquad
  \bigl(i<N\wedge A(i,q)\bigr)\to E_A(N,q).
\]
Consequently one may derive the expected biimplication after existential
elimination on the reverse direction.
\end{theorem}

\begin{proof}
Use the least-search construction above.  Boolean reflection converts
\(A(i,q)\) to \(z(i,q)=0\).  In the forward direction instantiate the
existential at \(\mu(N,q)\).  In the reverse direction retain the existential
witness locally and apply \code{search\_found}; the resulting formula does not
mention the witness, so existential elimination closes it.
\end{proof}

\subsection{Bounded universal verification}

A bounded universal is handled by searching for a counterexample.  Let
\(c(i,q)=b_A(i,q)\), interpreted so that \(c=0\) means failure.  Put
\[
  U_A(N,q)\;:\!\Longleftrightarrow\;\mu_c(N,q)=N.
\]
If \(i<N\) and \(A(i,q)\) failed, the reflected zero and
\code{search\_found} would imply \(\mu_c<N\), contradicting \(U_A\).  Conversely,
if every \(i<N\) satisfies \(A\), \(\mu_c<N\) would expose a counterexample by
\code{search\_spec}.  The reasoning is intuitionistic because equality and
order on natural-number terms are decidable in the reflected fragment.

For trace checking it is often cheaper to use a failure sum
\(\sum_{i<N}\mathsf{fail}(i,q)\) instead of a least counterexample.  The two
forms are interderivable for Boolean summands.  Search is preferable when the
counterexample index will be used; a sum is preferable when only global
validity matters.

\subsection{Fixed bounded syntax}

A useful Lean implementation has a host datatype with constructors
\[
  \mathsf{atom},\quad \wedge,\quad\vee,\quad\to,\quad
  \exists{<}t,\quad\forall{<}t.
\]
Each atom contains its reflected bit package.  The compiler recursively emits:
\begin{enumerate}[label=(\alph*)]
\item the ordinary formula used for human-facing statements;
\item a quantifier-free characteristic formula or verifier term;
\item derivations in both directions;
\item a proof of quantifier-freeness for every invariant later passed to
      \code{PRA.qfInd}.
\end{enumerate}
The recursion is a Lean recursion on the fixed host syntax.  Bounded
quantifiers invoke the search adapters above.  This precisely avoids the
impossible requirement of a total universal primitive-recursive evaluator for
all PR codes.

\subsection{Arrays and finite tuples}

The framework needs two coding modes.

\begin{description}
\item[Structural tuples.]  For a fixed number of fields, use nested pairing and
the checked projection equations.  This is best for query and state records.
\item[Bounded arrays.]  If every entry is below \(B\), code
\(a_0,\ldots,a_{N-1}\) by \(\sum_{i<N}a_i B^i\).  Digit extraction is quotient
and remainder by \(B^i\).  This is best for traces and strategy tables.
\end{description}

Mixed-radix coding is a derived optimization when field bounds differ.  SCR
requires a proof of every radix bound.  It never infers a code cutoff from a
statement that there are only finitely many mathematical objects.
